\documentclass[12pt]{article}
\usepackage[utf8]{inputenc}
\usepackage[spanish]{babel}
\usepackage{listings}
\usepackage{xcolor}
\usepackage{amsmath}
\usepackage{amssymb}
\usepackage{geometry}
\usepackage{fancyvrb}
\geometry{margin=1in}

% Configuración de colores
\definecolor{codegreen}{rgb}{0,0.6,0}
\definecolor{codegray}{rgb}{0.5,0.5,0.5}
\definecolor{codepurple}{rgb}{0.58,0,0.82}
\definecolor{backcolour}{rgb}{0.95,0.95,0.92}
\definecolor{outputcolor}{rgb}{0.1,0.1,0.4}

\lstdefinestyle{haskellstyle}{
    backgroundcolor=\color{backcolour},
    commentstyle=\color{codegreen},
    keywordstyle=\color{magenta},
    numberstyle=\tiny\color{codegray},
    stringstyle=\color{codepurple},
    basicstyle=\ttfamily\footnotesize,
    breakatwhitespace=false,
    breaklines=true,
    captionpos=b,
    keepspaces=true,
    numbers=left,
    numbersep=5pt,
    showspaces=false,
    showstringspaces=false,
    showtabs=false,
    tabsize=2
}

\lstset{style=haskellstyle}

\title{Demostraciones: Recursión, Autointerpretación y Metaprogramación\\
\large con el Combinador Y}
\author{}
\date{}

\begin{document}

\maketitle

\section*{Guión para Demostraciones en Vivo}

Este documento presenta el diálogo y las demostraciones para mostrar el funcionamiento del Combinador Y en tres contextos: recursión básica, autointerpretación y metaprogramación.

\newpage
\section{Demostración 1: Recursión con el Combinador Y}

\subsection*{Introducción}

\begin{quote}
\textbf{Presentador:} ``Vamos a comenzar demostrando cómo funciona la recursión usando el combinador Y. Lo especial aquí es que la función de evaluación \textbf{nunca se llama a sí misma por nombre}.''

\textbf{Presentador:} ``Ejecutemos el programa de test de expresiones aritméticas...''
\end{quote}

\subsection*{Ejecución del Programa}

\begin{verbatim}
$ ./test_expresiones
==========================================
TEST: EVALUACION DE EXPRESIONES ARITMETICAS
CON COMBINADOR Y
==========================================

TEST 1: Expresion simple
========================
Expresion: (+ 2 3)
Resultado: 5
\end{verbatim}

\subsection*{Explicación Paso a Paso}

\begin{quote}
\textbf{Presentador:} ``Veamos cómo se expande esta evaluación internamente...''
\end{quote}

\begin{verbatim}
Expansion paso a paso:
  evalExpr (+ 2 3)
  = (yFix evalMaker) (+ 2 3)
  = (evalMaker evalExpr) (+ 2 3)    [por punto fijo]
  = evalExpr 2 + evalExpr 3
  = 2 + 3
  = 5
\end{verbatim}

\begin{quote}
\textbf{Presentador:} ``Observen los pasos clave:''
\end{quote}

\textbf{Paso 1:} \texttt{evalExpr} se define como \texttt{yFix evalMaker}

\begin{equation*}
\texttt{evalExpr} = \texttt{yFix evalMaker}
\end{equation*}

\textbf{Paso 2:} Por la definición de \texttt{yFix}, esto se expande a:

\begin{equation*}
\texttt{yFix f} = \texttt{f (yFix f)}
\end{equation*}

Por lo tanto:
\begin{equation*}
\texttt{yFix evalMaker} = \texttt{evalMaker (yFix evalMaker)} = \texttt{evalMaker evalExpr}
\end{equation*}

\textbf{Paso 3:} Cuando \texttt{evalMaker} procesa \texttt{(+ 2 3)}, usa su parámetro \texttt{evalRec} (que es \texttt{evalExpr}) para evaluar las subexpresiones:

\begin{align*}
\texttt{evalMaker evalRec (+ 2 3)} &= \texttt{evalRec 2 + evalRec 3} \\
&= \texttt{evalExpr 2 + evalExpr 3} \\
&= 2 + 3 \\
&= 5
\end{align*}

\begin{quote}
\textbf{Presentador:} ``Lo crucial aquí es que \texttt{evalMaker} \textbf{nunca se llama a sí misma}. Usa el parámetro \texttt{evalRec} que el combinador Y le proporciona.''
\end{quote}

\newpage
\subsection*{Ejemplo Anidado}

\begin{verbatim}
TEST 2: Expresion anidada
==========================
Expresion: (* (+ 3 4) (- 10 2))
Resultado: 56
\end{verbatim}

\begin{quote}
\textbf{Presentador:} ``Con expresiones anidadas, la recursión se hace más evidente...''
\end{quote}

\begin{verbatim}
Expansion paso a paso:
  evalExpr (* (+ 3 4) (- 10 2))
  = (evalMaker evalExpr) (* (+ 3 4) (- 10 2))
  = evalExpr (+ 3 4) * evalExpr (- 10 2)
  = ((evalMaker evalExpr) (+ 3 4)) * ((evalMaker evalExpr) (- 10 2))
  = (evalExpr 3 + evalExpr 4) * (evalExpr 10 - evalExpr 2)
  = (3 + 4) * (10 - 2)
  = 7 * 8
  = 56
\end{verbatim}

\begin{quote}
\textbf{Presentador:} ``Cada llamada recursiva sigue el mismo patrón:''
\end{quote}

\begin{enumerate}
\item \texttt{evalExpr} expande a \texttt{evalMaker evalExpr}
\item \texttt{evalMaker} usa \texttt{evalExpr} (vía parámetro) para subexpresiones
\item El ciclo se repite hasta llegar a números (casos base)
\end{enumerate}

\subsection*{Observaciones Finales}

\begin{verbatim}
==========================================
OBSERVACIONES CLAVE:
==========================================
1. evalMaker NUNCA se llama a si misma por nombre
2. La recursion viene del combinador Y via yFix
3. Propiedad de punto fijo: evalExpr = evalMaker evalExpr
4. En cada paso recursivo, evalMaker usa 'evalRec'
5. 'evalRec' es en realidad 'evalExpr' (inyectado por Y)
==========================================
\end{verbatim}

\begin{quote}
\textbf{Presentador:} ``Esto demuestra que la recursión no necesita ser una primitiva del lenguaje. Puede implementarse usando solo funciones de orden superior.''
\end{quote}

\newpage
\section{Demostración 2: Autointerpretación}

\subsection*{Introducción}

\begin{quote}
\textbf{Presentador:} ``Ahora vamos a ver algo más profundo: la \textbf{autointerpretación}. Un intérprete que se usa a sí mismo para interpretar subexpresiones.''

\textbf{Presentador:} ``La diferencia con la demostración anterior es conceptual: ahora enfatizamos que el intérprete \textit{se pasa a sí mismo como argumento}.''
\end{quote}

\subsection*{Ejecución del Programa}

\begin{verbatim}
$ ./autointerprete
=== AUTOINTERPRETE CON COMBINADOR Y ===

Expresion 1: (+ 2 3)
Resultado: 5

Expresion 2: (* (+ 3 4) (- 10 2))
Resultado: 56

Expresion 3: (/ 20 (+ 2 3))
Resultado: 4
\end{verbatim}

\subsection*{¿Por Qué Esto Es Autointerpretación?}

\begin{quote}
\textbf{Presentador:} ``La clave está en esta propiedad:''
\end{quote}

\begin{equation*}
\texttt{interprete} = \texttt{interpreteMaker interprete}
\end{equation*}

\begin{quote}
\textbf{Presentador:} ``Esto significa que \texttt{interpreteMaker} recibe \texttt{interprete} (es decir, \textit{a sí mismo}) como argumento.''
\end{quote}

\textbf{Análisis del código:}

\begin{lstlisting}[language=Haskell, caption=Autointérprete]
interpreteMaker :: (ASAValues -> Int) -> (ASAValues -> Int)
interpreteMaker self expr = case expr of
    NumV n -> n
    AddV e1 e2 -> self e1 + self e2
    --            ^^^^ Aqui 'self' es en realidad 'interprete'

interprete :: ASAValues -> Int
interprete = yFix interpreteMaker
-- interprete = interpreteMaker interprete
\end{lstlisting}

\begin{quote}
\textbf{Presentador:} ``Cuando evaluamos \texttt{(+ 2 3)}:''
\end{quote}

\begin{align*}
&\texttt{interprete (+ 2 3)} \\
&= \texttt{(interpreteMaker interprete) (+ 2 3)} \\
&= \texttt{interprete 2 + interprete 3} \quad \text{[interprete usa interprete]} \\
&= 2 + 3 = 5
\end{align*}

\begin{quote}
\textbf{Presentador:} ``El intérprete literalmente se usa a sí mismo. Esta es la definición de autointerpretación.''
\end{quote}

\subsection*{Diferencia con Recursión Tradicional}

\begin{center}
\begin{tabular}{|l|l|}
\hline
\textbf{Recursión Tradicional} & \textbf{Autointerpretación con Y} \\
\hline
\texttt{f x = ... f ...} & \texttt{fMaker self x = ... self ...} \\
Auto-referencia por nombre & Recibe referencia como parámetro \\
Ciclo implícito en el lenguaje & Ciclo explícito vía punto fijo \\
\texttt{f} está predefinida & \texttt{f = yFix fMaker} \\
\hline
\end{tabular}
\end{center}

\begin{quote}
\textbf{Presentador:} ``Sin el combinador Y, la autointerpretación sería imposible en cálculo lambda puro, porque no habría forma de que una función se reciba a sí misma como argumento.''
\end{quote}

\subsection*{Mensaje Final}

\begin{verbatim}
Propiedad de autointerpretacion:
  interprete = interpreteMaker interprete
  El interprete se pasa a si mismo como argumento
  Esto es AUTOINTERPRETACION gracias al combinador Y
\end{verbatim}

\newpage
\subsection*{Tests Formales de Autointerpretación}

\begin{quote}
\textbf{Presentador:} ``Para verificar que nuestra autointerpretación funciona correctamente, ejecutemos los tests formales...''
\end{quote}

\begin{verbatim}
$ ./test_autointerprete
==========================================
TESTS FORMALES: AUTOINTERPRETACION
==========================================

Test 1 - Numeros simples: Pass
Test 2 - Suma simple: Pass
Test 3 - Resta: Pass
Test 4 - Multiplicacion: Pass
Test 5 - Division: Pass
Test 6 - Expresion anidada: Pass
Test 7 - Expresion compleja: Pass
Test 8 - Division anidada: Pass
Test 9 - Punto fijo (interprete = interpreteMaker interprete): Pass
Test 10 - Expresion muy anidada: Pass

==========================================
RESULTADO: 10/10 tests pasados
EXITO: Todos los tests de autointerpretacion pasaron!
\end{verbatim}

\begin{quote}
\textbf{Presentador:} ``El Test 9 es particularmente importante. Veamos por qué...''
\end{quote}

\textbf{Test 9: Verificación de la Propiedad de Punto Fijo}

\begin{lstlisting}[language=Haskell]
-- Test 9: Verificar autointerpretacion
let expr9 = AddV (NumV 10) (NumV 20)
let result9a = interprete expr9
let result9b = interpreteMaker interprete expr9
let test9 = runTest "Test 9: Propiedad de punto fijo"
                    (result9a == result9b)
\end{lstlisting}

\begin{quote}
\textbf{Presentador:} ``¿Por qué este test verifica autointerpretación?''
\end{quote}

\textbf{Explicación:}

\begin{enumerate}
\item \textbf{Calculamos de dos formas}:
\begin{itemize}
\item \texttt{result9a = interprete expr9} (usando directamente el intérprete)
\item \texttt{result9b = interpreteMaker interprete expr9} (aplicando interpreteMaker a interprete)
\end{itemize}

\item \textbf{Si son iguales}, entonces se cumple: $\texttt{interprete} = \texttt{interpreteMaker interprete}$

\item \textbf{Esto prueba que} \texttt{interpreteMaker} recibe \texttt{interprete} (a sí mismo) como argumento

\item \textbf{Por lo tanto}, el intérprete se está usando a sí mismo para evaluar subexpresiones = \textbf{AUTOINTERPRETACIÓN}
\end{enumerate}

\begin{quote}
\textbf{Presentador:} ``Sin el combinador Y, esta propiedad sería imposible de lograr en cálculo lambda puro, porque no habría forma de crear esta auto-referencia sin nombrar la función recursivamente.''
\end{quote}

\newpage
\section{Demostración 3: Metaprogramación}

\subsection*{Introducción}

\begin{quote}
\textbf{Presentador:} ``Finalmente, veamos cómo la autointerpretación habilita la \textbf{metaprogramación}: la capacidad del código para manipularse a sí mismo.''

\textbf{Presentador:} ``Demostraremos dos tipos de metaprogramación:''
\end{quote}

\begin{enumerate}
\item Inspección de código - El programa examina su estructura
\item Optimización de código - El programa se transforma a sí mismo
\end{enumerate}

\subsection*{Ejecución del Programa}

\begin{verbatim}
$ ./metaprogramacion
==========================================
METAPROGRAMACION CON COMBINADOR Y
==========================================
\end{verbatim}

\subsection*{Parte 1: Inspección de Código}

\begin{quote}
\textbf{Presentador:} ``Primero, el código puede \textit{inspeccionarse} a sí mismo para detectar características peligrosas...''
\end{quote}

\begin{verbatim}
1. INSPECCION DE CODIGO
====================
Expresion: (+ 5 3)
Contiene division? False
Evaluacion segura: 8

Expresion: (/ 10 2)
Contiene division? True
ADVERTENCIA: Contiene division
\end{verbatim}

\textbf{Explicación:}

\begin{lstlisting}[language=Haskell, caption=Inspección de código]
contieneDivision :: ASAValues -> Bool
contieneDivision expr = case expr of
    NumV _ -> False
    DivV _ _ -> True  -- Encontro una division!
    AddV e1 e2 -> contieneDivision e1 || contieneDivision e2
    SubV e1 e2 -> contieneDivision e1 || contieneDivision e2
    MultV e1 e2 -> contieneDivision e1 || contieneDivision e2
\end{lstlisting}

\begin{quote}
\textbf{Presentador:} ``Esta función \textit{analiza} el código antes de ejecutarlo. El programa examina su propia estructura.''

\textbf{Presentador:} ``¿Por qué esto es metaprogramación? Porque el código (tipo \texttt{ASAValues}) se trata como \textit{datos} que podemos inspeccionar.''
\end{quote}

\newpage
\subsection*{Parte 2: Optimización de Código}

\begin{quote}
\textbf{Presentador:} ``Ahora veamos cómo el código puede \textit{transformarse} a sí mismo...''
\end{quote}

\begin{verbatim}
2. OPTIMIZACION DE CODIGO
======================
Expresion original: (+ (* 0 999) 5)
Expresion optimizada: (+ 0 5)
Resultado sin optimizar: 5
Resultado optimizado: 5
\end{verbatim}

\textbf{Explicación del proceso:}

\begin{enumerate}
\item \textbf{Expresión original:} \texttt{(+ (* 0 999) 5)}
\item \textbf{Optimización detecta:} Multiplicación por 0
\item \textbf{Regla aplicada:} \texttt{(* 0 x) $\rightarrow$ 0}
\item \textbf{Resultado:} \texttt{(+ 0 5)}
\item \textbf{Segunda optimización:} Suma con 0
\item \textbf{Regla aplicada:} \texttt{(+ 0 x) $\rightarrow$ x}
\item \textbf{Resultado final:} \texttt{5}
\end{enumerate}

\begin{quote}
\textbf{Presentador:} ``El programa se reescribe a sí mismo antes de ejecutarse, eliminando cálculos innecesarios.''
\end{quote}

\textbf{Más ejemplos:}

\begin{verbatim}
Expresion original: (* (+ 0 7) (+ 3 0))
Expresion optimizada: (* 7 3)
Resultado sin optimizar: 21
Resultado optimizado: 21

Expresion original: (+ (* 1 10) (* 0 100))
Expresion optimizada: (+ 10 0)
Resultado sin optimizar: 10
Resultado optimizado: 10
\end{verbatim}

\begin{quote}
\textbf{Presentador:} ``Observen cómo se simplifican las expresiones:''
\end{quote}

\begin{itemize}
\item \texttt{(+ 0 x)} se simplifica a \texttt{x}
\item \texttt{(* 1 x)} se simplifica a \texttt{x}
\item \texttt{(* 0 x)} se simplifica a \texttt{0}
\end{itemize}

\subsection*{Código de Optimización}

\begin{lstlisting}[language=Haskell, caption=Transformación de código]
optimizar :: ASAValues -> ASAValues
optimizar expr = case expr of
    -- Eliminar sumas con cero
    AddV (NumV 0) e -> optimizar e
    AddV e (NumV 0) -> optimizar e

    -- Eliminar multiplicaciones por 1
    MultV (NumV 1) e -> optimizar e
    MultV e (NumV 1) -> optimizar e

    -- Multiplicacion por 0 = 0
    MultV (NumV 0) _ -> NumV 0
    MultV _ (NumV 0) -> NumV 0

    -- Recursivamente optimizar subexpresiones
    AddV e1 e2 -> AddV (optimizar e1) (optimizar e2)
    ...
\end{lstlisting}

\begin{quote}
\textbf{Presentador:} ``Esta función \textit{transforma} el código. Toma código como entrada y produce código como salida.''
\end{quote}

\subsection*{¿Por Qué el Combinador Y Es Esencial?}

\begin{quote}
\textbf{Presentador:} ``Estas capacidades de metaprogramación son posibles porque:''
\end{quote}

\begin{enumerate}
\item \textbf{Autointerpretación:} El combinador Y permite que el lenguaje se interprete a sí mismo
\item \textbf{Código = Datos:} Las expresiones son valores de tipo \texttt{ASAValues}
\item \textbf{Recursión ilimitada:} El combinador Y permite procesar estructuras arbitrariamente anidadas
\item \textbf{Sin primitivas:} Todo funciona en cálculo lambda puro
\end{enumerate}

\newpage
\subsection*{Tests Formales de Metaprogramación}

\begin{quote}
\textbf{Presentador:} ``Ahora verifiquemos formalmente que nuestras capacidades de metaprogramación funcionan correctamente...''
\end{quote}

\begin{verbatim}
$ ./test_metaprogramacion
==========================================
TESTS FORMALES: METAPROGRAMACION
==========================================

PARTE 1: TESTS DE INSPECCION DE CODIGO
========================================
Test 1 - Detectar division presente: Pass
Test 2 - Detectar ausencia de division: Pass
Test 3 - Detectar division anidada: Pass
Test 4 - Evaluacion segura sin division: Pass
Test 5 - Advertencia con division: Pass

PARTE 2: TESTS DE OPTIMIZACION DE CODIGO
==========================================
Test 6 - Optimizar suma con cero: Pass
Test 7 - Optimizar multiplicacion por 1: Pass
Test 8 - Optimizar multiplicacion por 0: Pass
Test 9 - Optimizacion anidada: Pass
Test 10 - Resultado preservado tras optimizacion: Pass
Test 11 - Optimizacion compleja: Pass
Test 12 - Preservacion resultado complejo: Pass

==========================================
RESULTADO: 12/12 tests pasados
EXITO: Todos los tests de metaprogramacion pasaron!
\end{verbatim}

\begin{quote}
\textbf{Presentador:} ``Estos tests demuestran dos cosas fundamentales...''
\end{quote}

\textbf{1. ¿Por qué la inspección de código es metaprogramación?}

\begin{lstlisting}[language=Haskell]
-- Test: Detectar division presente
let test1 = runTest "Test 1: Detectar division en (/ 10 2)"
                    (contieneDivision (DivV (NumV 10) (NumV 2)) == True)
\end{lstlisting}

\textbf{Explicación:}
\begin{enumerate}
\item La función \texttt{contieneDivision} recibe \textbf{código} (expresión ASA) como entrada
\item \textbf{Analiza} la estructura del código sin ejecutarlo
\item Retorna información \textbf{sobre el código} (si contiene división)
\item El código está \textbf{inspeccionando código} = \textbf{METAPROGRAMACIÓN}
\end{enumerate}

\begin{quote}
\textbf{Presentador:} ``Esto es posible porque el combinador Y nos permitió construir un intérprete que puede razonar sobre su propio código.''
\end{quote}

\textbf{2. ¿Por qué la optimización es metaprogramación?}

\begin{lstlisting}[language=Haskell]
-- Test: Optimizacion anidada
let expr9 = AddV (MultV (NumV 0) (NumV 999)) (NumV 5)
let test9 = runTest "Test 9: Optimizar (+ (* 0 999) 5)"
                    (optimizar expr9 == NumV 5)
\end{lstlisting}

\textbf{Explicación paso a paso:}
\begin{enumerate}
\item \textbf{Entrada}: Código \texttt{(+ (* 0 999) 5)}
\item \textbf{Función optimizar}:
\begin{itemize}
\item Inspecciona: detecta \texttt{(* 0 999)}
\item Aplica regla: \texttt{(* 0 x) $\rightarrow$ 0}
\item Genera nuevo código: \texttt{(+ 0 5)}
\item Inspecciona: detecta \texttt{(+ 0 5)}
\item Aplica regla: \texttt{(+ 0 x) $\rightarrow$ x}
\item Genera código final: \texttt{5}
\end{itemize}
\item \textbf{Salida}: Código transformado \texttt{5}
\item El código está \textbf{transformando código} = \textbf{METAPROGRAMACIÓN}
\end{enumerate}

\textbf{3. Verificación de Preservación de Resultados}

\begin{lstlisting}[language=Haskell]
-- Test: Verificar que optimizacion preserva resultado
let expr10 = MultV (AddV (NumV 0) (NumV 7)) (AddV (NumV 3) (NumV 0))
let test10 = runTest "Test 10: Resultado preservado"
                     (interprete expr10 == evalOptimizado expr10)
\end{lstlisting}

\begin{quote}
\textbf{Presentador:} ``Este test es crucial: verifica que las transformaciones de código \textit{preservan el significado}. El código optimizado debe dar el mismo resultado que el original.''
\end{quote}

\textbf{¿Cómo funciona la preservación?}
\begin{enumerate}
\item \texttt{interprete expr10} evalúa el código original $\rightarrow$ resultado 21
\item \texttt{evalOptimizado expr10 = interprete (optimizar expr10)}
\begin{itemize}
\item Primero optimiza: \texttt{(* (+ 0 7) (+ 3 0)) $\rightarrow$ (* 7 3)}
\item Luego evalúa: \texttt{(* 7 3) $\rightarrow$ 21}
\end{itemize}
\item Si ambos dan 21, la transformación es \textbf{correcta}
\item Todos los tests verifican esta propiedad $\rightarrow$ \textbf{12/12 tests pasados}
\end{enumerate}

\subsection*{Conclusiones}

\begin{verbatim}
==========================================
CONCLUSIONES:
- El codigo puede INSPECCIONARSE a si mismo
- El codigo puede TRANSFORMARSE a si mismo
- Esto es METAPROGRAMACION habilitada por el combinador Y
==========================================
\end{verbatim}

\begin{quote}
\textbf{Presentador:} ``La metaprogramación que acabamos de ver es la base de:''
\end{quote}

\begin{itemize}
\item Compiladores optimizadores
\item Sistemas de macros (como en Lisp/Scheme)
\item Compilación JIT (Just-In-Time)
\item Evaluación parcial
\item Análisis estático de código
\item Reflection en lenguajes modernos
\end{itemize}

\begin{quote}
\textbf{Presentador:} ``Todo esto es posible gracias al simple pero poderoso combinador Y, que permite que un lenguaje se interprete y se transforme a sí mismo.''
\end{quote}

\newpage
\section{Resumen: Cadena de Implicaciones}

\begin{quote}
\textbf{Presentador:} ``Para finalizar, veamos cómo todo se conecta:''
\end{quote}

\begin{align*}
\text{Combinador Y} &\Rightarrow \text{Recursión sin auto-referencia} \\
&\Rightarrow \text{Autointerpretación posible} \\
&\Rightarrow \text{Código como datos manipulables} \\
&\Rightarrow \text{Metaprogramación}
\end{align*}

\subsection*{Tabla Comparativa Final}

\begin{center}
\begin{tabular}{|l|p{8cm}|}
\hline
\textbf{Concepto} & \textbf{Qué Demuestra} \\
\hline
\textbf{Recursión} & \texttt{evalMaker} nunca se llama a sí misma. La recursión viene del combinador Y. \\
\hline
\textbf{Autointerpretación} & El intérprete se pasa a sí mismo como argumento: \texttt{interprete = interpreteMaker interprete} \\
\hline
\textbf{Metaprogramación} & El código se inspecciona y transforma a sí mismo antes de ejecutarse. \\
\hline
\end{tabular}
\end{center}

\subsection*{Mensaje Final}

\begin{quote}
\textbf{Presentador:} ``Hemos demostrado que:''
\end{quote}

\begin{enumerate}
\item La recursión no requiere auto-referencia explícita
\item Un lenguaje puede interpretarse a sí mismo gracias al combinador Y
\item La autointerpretación habilita metaprogramación
\item Estas capacidades funcionan en cálculo lambda puro
\end{enumerate}

\begin{quote}
\textbf{Presentador:} ``El combinador Y no es solo una curiosidad teórica. Es el fundamento de características avanzadas en lenguajes modernos: optimización de compiladores, reflexión, macros, y metaprogramación.''

\textbf{Presentador:} ``Todo esto emerge de una simple ecuación de punto fijo: \texttt{f = yFix fMaker}.''
\end{quote}

\end{document}
